\section{Conclusion}
\label{sec:conclusion}

The certified computation determines \(\MB(n)\) for every nonempty Barnette
census from order 8 through order 36, including all ties and the unique
order-36 value 49.  Its proof objects are explicit: upper certificates are
Hamiltonian-cycle families, and lower certificates are either packings or
complete-universe exhaustive exclusions checked by a standalone verifier.

That computation exposed a structural double-ladder recurrence.  A rail/rung
recurrence now proves that \(D(a,b)\), for odd \(a,b\ge3\), has one all-rails
cycle, \(ab\) turn cycles, and four connector exceptions.  The resulting
coordinate system supports two matching symbolic constructions.  Exception
and boundary singletons together with an interior domino tiling form a packing
of \(((a+2)(b+2)-1)/2\) requirements.  The four exceptions, all boundary
turns, and one interior parity class form a separating family of the same
size; its rail, rung, and connector signatures are pairwise incomparable.
Thus the exact formula
\[
 \hsep(D(a,b))=\frac{(a+2)(b+2)-1}{2}
\]
holds throughout this infinite family.

The three prediction-locked graphs beyond order 36 provide prospective finite
confirmation of the cycle and edge-separation formulas, but no census-level
claim at orders 40, 44, or 48.  Likewise, the exact family theorem does not
prove that balanced double ladders remain extremal among all Barnette graphs.
That continuation, and a structural explanation of the tied boundary-state
branch at orders congruent to two modulo four, remain open.

Before circulation as a novelty claim, the manuscript still requires a
systematic literature review of related separation systems, cycle-separation
parameters, Barnette censuses, and graph-expansion operations.  Independent
expert review of the symbolic case analysis and the direct 3-connectivity
argument also remains desirable.
